\documentclass[12pt,a4paper]{article}

\usepackage{fontawesome5}
\providecommand{\faPencilAlt}{\faIcon{pencil-alt}}

\usepackage{styles/couleurs}
\usepackage{styles/mise_en_page}
\usepackage{styles/boites}
\usepackage{styles/maths}
\usepackage{styles/figures}

\usepackage{enumitem}
\usepackage{array}
\usepackage{tabularx}
\usepackage{amsmath}
\usepackage{amssymb}
\usepackage{tikz}
\usepackage{xparse}

\usetikzlibrary{arrows.meta,calc}

\setlist[itemize]{leftmargin=1.4em,itemsep=0.15em,topsep=0.2em}
\setlist[enumerate]{leftmargin=*,itemsep=0.25em,topsep=0.25em}

\newcommand{\lignepointillee}{%
  \noindent\dotfill\par
  \vspace{0.45em}
}

\newcommand{\lignesreponse}[1]{%
  \begin{tcolorbox}[
    enhanced,
    colback=white,
    colframe=NeutreC!40,
    boxrule=0.5pt,
    arc=1.5mm,
    left=3mm,
    right=3mm,
    top=2mm,
    bottom=1mm,
    before skip=2pt,
    after skip=2pt
  ]
    \foreach \i in {1,...,#1}{%
      \lignepointillee
    }
  \end{tcolorbox}
}

\newcommand{\titreseq}[1]{%
\begin{center}
  {\sffamily\LARGE\bfseries\color{TitreC}Séquence 3 -- Nombres relatifs}\\[0.2em]
  {\sffamily\large\bfseries\color{SavoirC}#1}
\end{center}
}

\begin{document}

\titreseq{Activité 7 -- Formule Celsius-Fahrenheit}

\remarque{Cette activité sert de synthèse : elle mobilise les produits, les sommes, les différences, les quotients et l'interprétation du résultat.}

En France, on mesure la température en degrés Celsius. Dans d'autres pays, elle se mesure en degrés Fahrenheit.

La formule est :

\[
t_F=1{,}8\times t_C+32.
\]

\section*{Partie A -- Calculer une température en degrés Fahrenheit}

Calcule la température en degrés Fahrenheit correspondant à :

\[
10^\circ\mathrm{C}
\qquad
0^\circ\mathrm{C}
\qquad
-10^\circ\mathrm{C}.
\]

\lignesreponse{7}

\section*{Partie B -- Retrouver une température en degrés Celsius}

On sait que :

\[
t_F=59.
\]

Complète le programme inverse.

\[
59-32=\ldots
\]

\[
\ldots\div1{,}8=\ldots
\]

Donc :

\[
59^\circ\mathrm{F}=\ldots^\circ\mathrm{C}.
\]

\lignesreponse{4}

\section*{Partie C -- Écrire une expression unique}

Écris une seule expression permettant de retrouver la température en degrés Celsius correspondant à \(59^\circ\mathrm{F}\).

\[
\dotfill
\]

Calcule-la.

\lignesreponse{3}

\section*{Partie D -- Calculatrice raisonnée}

Une journée d'hiver, une ville américaine annonce \(-4^\circ\mathrm{F}\).

\begin{enumerate}
  \item Avant de calculer, penses-tu que la température en Celsius sera positive ou négative ?
  \item Calcule la température en Celsius avec une expression unique.
  \item Interprète le résultat.
\end{enumerate}

\lignesreponse{6}

\end{document}
